\begin{table}[htbp]
\centering
\resizebox{\columnwidth}{!}{%
\begin{tabular}{l|rr|rr|rr}
\hline
\multirow{2}{*}{\textbf{领域}} & \multicolumn{2}{c|}{\textbf{任务完成}} & \multicolumn{2}{c|}{\textbf{完整性约束}} & \multicolumn{2}{c}{\textbf{策略合规}} \\
& \textbf{样本数} & \textbf{分数} & \textbf{样本数} & \textbf{分数} & \textbf{样本数} & \textbf{分数} \\
\hline
Teams    & 100 & 82.3 & 20 & 90.0 & 8   & 81.3 \\
CSM      & 183 & 58.3 & 30 & 70.0 & 120 & 45.5 \\
Email    & 103 & 70.9 & 17 & 76.5 & 12  & 79.2 \\
ITSM     & 176 & 53.7 & 30 & 50.6 & 70  & 30.0 \\
Calendar & 104 & 71.5 & 8  & 62.5 & 19  & 76.1 \\
HR       & 182 & 63.2 & 47 & 56.0 & 94  & 50.4 \\
Drive    & 105 & 75.1 & 26 & 87.8 & 6   & 66.7 \\
Hybrid   & 152 & 60.0 & 23 & 72.5 & 18  & 61.1 \\
\hline
\end{tabular}
}
\caption{按类别统计的跨领域验证器通过率。模型在完整性约束（系统状态有效性）上得分最高，其次为任务完成；策略合规表现最低。各类别样本数不同；CSM、ITSM 和 HR 等负担更重的领域更强调完整性约束和策略合规。}
\label{tab:verifier_analysis}
\end{table}
